\subsection{Sprint 0 --- Fase de Planejamento Estratégico}

O projeto teve início no dia 02/03/2026, com uma reunião curta para nos conhecermos e
lermos o Termo de Abertura. A impressão inicial foi de que se tratava de um projeto
simples, sem grandes impeditivos. Na reunião seguinte, aprofundamos o entendimento do
projeto e preparamos as perguntas para quando os clientes viessem apresentar.

Na semana seguinte, os clientes apresentaram suas ideias com muito entusiasmo e diversas
propostas. A ausência de projeção de slides dificultou o entendimento, pois havia muitas
ideias a tratar. Os colegas de \ac{ages} IV perceberam que o escopo estava se expandindo
demais e relembraram que o projeto era um \ac{mvp} --- uma Prova de Conceito --- e que o
Termo de Abertura deveria ser seguido.

Os clientes chegaram a sugerir a exclusão do site existente e uma integração completa ao
novo sistema. A orientadora trouxe o contexto de que a aplicação web seria algo separado,
sugerindo manter o site atual. A página de cursos FATLAB funcionaria como ponto de
entrada para o novo sistema, onde seria possível gerir os alunos e gerar dados para
indicadores futuros. Também construiríamos um acesso para empresas disponibilizarem
vagas, com lógica de negócio baseada em remuneração por indicação de alunos ---
garantindo entrada financeira para manter a instituição. A reunião ocupou toda a aula,
impossibilitando o alinhamento interno naquele dia.

Nesse período, fiz muitas anotações e explorei o site da instituição. Criei rascunhos de
telas usando a \ac{ia} \textit{Stitch}, do Google, mas as imagens geradas não refletiam o que
eu imaginava. Utilizei um plugin para importar ao Figma, onde desenvolvi visualizações
para administração e alunos que pareciam atender melhor às necessidades dos clientes.

Na quarta-feira seguinte, já estava com as ideias prototipadas no Figma. O colega
Guilherme mencionou uma \ac{ia} capaz de gerar sites funcionais a partir de modelos visuais.
Testei a ferramenta com as telas que havia criado e o resultado foi um site funcional com
o mesmo padrão visual dos protótipos.

Em aula, desenvolvemos mockups de baixo nível no quadro e definimos a divisão de
tarefas para o Figma. O Andrei (\ac{ages} III) criou o projeto no Figma e definiu o padrão
visual (cores, tipografias, botões e ícones). Fiquei responsável, junto com Lucca e
Vinicius, pelas telas de \textbf{Cadastro de Aluno}, \textbf{Cadastro de Empresa},
\textbf{Edição de Cadastro de Aluno}, \textbf{Edição de Cadastro de Empresa} e
\textbf{Dashboard Administrativo}.

Durante duas semanas, mantive foco quase exclusivo nas telas de cadastro, que sofreram
diversas iterações conforme surgiam ajustes no design. No dia 23/03, soube que o
designer Andrei havia comentado com os \ac{ages} IV que eu estaria fazendo as telas de
forma errada, sem jamais ter me comunicado diretamente. Como não tenho resistência a
críticas construtivas, fui até ele e sugeri que marcasse um momento para me mostrar as
correções. Ele reconheceu que o problema era coletivo, não individual.

A partir daí, o time entrou em modo de finalização das telas. Fiz as correções e, no
dia da apresentação, corrigi desalinhamentos de prototipagem em tempo real. A
orientadora identificou problemas críticos durante a apresentação, mas conseguimos
reorganizar rapidamente. Os \ac{ages} IV apresentaram de forma dinâmica e sincronizada, e
a cliente saiu satisfeita com a proposta.

Como \ac{ages} II, criei um grupo para integrar os colegas e facilitar os alinhamentos
sobre o banco de dados. Organizei uma reunião no dia 21/03, com duração de 1h30.
Todos participaram, exceto um colega que não opinou mas esteve presente. Definimos
a estrutura do banco e migramos o template do MIRO para o \textit{dbdiagram}. O
Tarciso pediu minha opinião sobre a tela de home do aluno --- ele achava que não
combinava com o restante do projeto. Pesquisei exemplos e compartilhei com ele,
envolvendo também a Isa (\ac{ages} IV), que ajudou a melhorar a tela. O \ac{uml} foi
aprovado com ressalvas, pois mudanças futuras seriam necessárias.

Na semana seguinte, o Heitor pediu mudanças nos atributos e traduções. Realizamos
nova reunião no domingo, dia 29/03, às 21h, com Lucca, Pedro, Tarciso e eu. Fizemos
alterações e pensamos em novas lógicas surgidas na reunião com o cliente.

Por fim, realizamos a retrospectiva da Sprint 0, que foi muito animada. Utilizamos um
sistema de avaliação inspirado no cinema, com dois pontos centrais:

\begin{enumerate}
  \item Aspectos positivos da sprint e destaque para quem se sobressaiu;
  \item Aspectos a melhorar, com destaque para comportamentos a evitar, como atrasos
    ou falta de participação.
\end{enumerate}

Recebi elogios de colegas que eu não imaginava terem notado meu esforço. Parabenizei
o Guilherme (\ac{ages} IV), que está atuando muito bem como arquiteto do projeto. Houve
uma discussão porque um colega ficou furioso por ter sido citado como exemplo negativo,
mas fora isso, tudo correu bem. Estamos alinhados e muito ansiosos pelas próximas etapas.
